\section{The Agent-Usability Framework}
\label{sec:auf}

An industrial agent receives a goal, resolves it against a plant it has not seen, decides
which evidence it needs, and proposes an answer or action that must be checkable before
dispatch. AUF asks what each source of engineering semantics contributes to that cycle.
The ten capabilities were derived by tracing the minimum evidence obligations of a
bounded supervisory cycle (resolve, interpret, locate, decide, check, trust, retrieve)
and the two cross-cutting feasibility conditions of a deployment (context budget and
ecosystem). Table~\ref{tab:rubric} states for each capability the operational requirement,
the agent failure it prevents, the scoring anchors and representative sources. The
derivation is tested against the coded literature as a saturation check: any recurring
requirement that cannot be mapped to U1--U10 is reported as a framework gap.

\begin{table*}[t]
\centering\footnotesize
\setlength{\tabcolsep}{3pt}
\caption{AUF capabilities: requirement, failure mode prevented and rubric anchors (0 absent;
1 incidental; 2 supported; 3 first class). Scores reflect what a specification
standardises by default, not optional extensions.}
\label{tab:rubric}
\begin{tabular}{@{}L{1.9cm}L{3.2cm}L{3.0cm}L{2.6cm}L{4.4cm}L{1.6cm}@{}}
\toprule
Capability & Requirement & Failure prevented & 1 incidental & 3 first class & Sources \\
\midrule
U1 Identity & resolve phrase to the right concept and instance & wrong equipment; homograph conflation & identifiers unique per vendor or file & globally unique, versioned, resolvable, synonym handling & \cite{iso29002,eclass,iec61360} \\
U2 Types and access & know datatypes, ranges, access rights & malformed call; write to read-only point & datatype only & datatypes, ranges, access rights, enumerations, run-time introspection & \cite{iec62541,patil2024gorilla} \\
U3 Quantities & bind values to quantity kind and unit & dimension or scale error & unit as display string & quantity kind, dimension vector, affine conversion & \cite{qudt} \\
U4 Relations & containment, topology, typed relations & wrong scope; missed impact & containment only & typed relations with transitivity and inverses & \cite{iec62264,balaji2016brick} \\
U5 Behaviour & operations, state machines, legal commands & command illegal in current state & operations named, untyped & typed operations plus state model with legality conditions & \cite{isa88,vdi2658} \\
U6 Constraints & checkable guards and envelopes & unsafe or out-of-envelope proposal & range or type checks & executable constraint language with reports & \cite{knublauch2017shacl,w3c2018odrl} \\
U7 Provenance and time & timestamped, attributed evidence & stale or untraceable evidence & one global timestamp & per-value time, source, agent and derivation & \cite{w3c2013prov,haller2019ssn} \\
U8 Query pushdown & filter, join, aggregate outside the model & context stuffing; unanswerable aggregates & filtered fetch & query language with joins, aggregation, path traversal & \cite{w3c2013sparql,edge2024graphrag} \\
U9 Context economy & faithful minimal evidence slice & budget exhaustion before evidence & verbose, partial fetch & compact and projectable to a minimal faithful slice & \cite{lewis2020rag,sarmah2024hybridrag} \\
U10 Ecosystem & implementations, converters, governance & deployment never happens & one implementation & open implementations, converters, reference data, active governance & \cite{idta2023aas,mcp} \\
\bottomrule
\end{tabular}
\end{table*}

Three design rules make AUF a repeatable instrument rather than an opinion table. First,
capabilities are scored separately, because strength in one does not compensate for
another: a \texttt{Double} in bar and a \texttt{Double} in psi have the same type, and a
complete address space need not state the order in which commands are legal. Second,
technologies are scored on their default specification, since that is what an agent
receives unless a deployment publishes more; deployment profiles are scored separately
when agent systems are coded. Third, U1--U8 are evidence obligations whereas U9--U10 are
feasibility conditions; AUF does not cover model training, cybersecurity or human
factors, and it does not prescribe a serialisation.
